% Source: ME16_v1.html
\documentclass[11pt]{article}
\usepackage[margin=1in]{geometry}
\usepackage{amsmath}
\usepackage{amssymb}
\usepackage{siunitx}
\usepackage{enumitem}

\title{ME16: Sound}
\author{}
\date{}

\begin{document}
\maketitle

\section*{ME16: Sound}

\begin{enumerate}[leftmargin=*,label=\textbf{Question \arabic*.}]
\item \hfill \textit{/ 10 pts}

A sound wave in an unknown fluid medium takes 14.6 ms to travel 20.4 m. What is the bulk modulus of the fluid in GPa, if the density is 778 kg/m$^{3}$? (1 GPa = 10$^{9}$ Pa)



\emph{B} = \rule{2cm}{0.4pt} GPa

\textit{Correct answer: } 1.5189

\item \hfill \textit{/ 10 pts}

A mine explodes off shore north of Cuba, creating a sound wave under water. How long does it take for the sound to register on Coast Guard sonar, 278 km away?

Density and bulk modulus of sea water: 1025 kg/m$^{3}$ and 2.34 x 10$^{9}$ Pa.



\emph{t} = \rule{2cm}{0.4pt} s

\textit{Correct answer: } 183.9919

\item \hfill \textit{/ 10 pts}

A siren is designed to create a decibel level of 123 dB at a distance of 11 m. Calculate the sound intensity in W/m$^{2}$at the same distance if there are 2 sirens all sounding at that same location.

\textit{Correct answer: } 3.9905

\item \hfill \textit{/ 10 pts}

A pneumatic drill creates a decibel level of 112 dB at a distance of 10 m. Calculate the \textbf{sound intensity level }in dBat a distance of 37 m.

\textit{Correct answer: } 100.636

\item \hfill \textit{/ 10 pts}

A cicada's mating call reaches 105 dB at a distance of 1.00 m from the insect. Calculate the decibel level at a distance of 25 m from 61 cicadas each producing the same sound intensity.

\textit{Correct answer: } 94.8945

\item \hfill \textit{/ 10 pts}

Consider three tuning forks, labeled A, B, and C. Forks A and B have frequencies of 440 Hz and 446 Hz, respectively. The beat frequency between forks A and C is 2.00 Hz while B and C have a beat frequency of 4.00 Hz. Calculate the frequency of fork C.

\begin{itemize}
  \item \textbf{[correct]} 442 Hz
  \item 443 Hz
  \item 444 Hz
  \item 448 Hz
  \item 438 Hz
\end{itemize}

\item \hfill \textit{/ 10 pts}

Speaker 1 is at \emph{x} = 0 and speaker 2 is to the right and facing it at \emph{x} = 7 m. The speakers are connected to the same amplifier, which plays a continuous tone at 347 Hz. Find the \emph{x}-coordinate in meters of the first point of constructive interference that is right of center, if the temperature in the room is 23°C.

Take the speed of sound to be $v=(331\;\mathrm{m/s})\sqrt{\frac{T}{273\;\mathrm{K}}}$ .



\emph{x} = \rule{2cm}{0.4pt} m

\textit{Correct answer: } 3.9966

\item \hfill \textit{/ 10 pts}

Speaker 1 is located at (\emph{x},\emph{y}) = (0, 2) m and speaker 2 is at (\emph{x},\emph{y}) = (0,0) m. They play the same frequency sound in the same phase. What is the minimum frequency in hertz that would result in constructive interference at the point (\emph{x},\emph{y}) = (5,0) m? Speed of sound: 344 m/s.



\emph{f} = \rule{2cm}{0.4pt} Hz

\textit{Correct answer: } 893.1242

\item \hfill \textit{/ 10 pts}

A pipe open at both ends has a length of 0.8 m. If the temperature of the room is 14°C, calculate the frequency in Hz of the \emph{n} = 2 harmonic.

Take the speed of sound to be $v=(331\;\mathrm{m/s})\sqrt{\frac{T}{273\;\mathrm{K}}}$ .

\textit{Correct answer: } 424.2263

\item \hfill \textit{/ 10 pts}

A pipe closed at one end has a length of 0.36 m. If the temperature of the room is 281 K, calculate the first overtone in Hz.

Take the speed of sound to be $v=(331\;\mathrm{m/s})\sqrt{\frac{T}{273\;\mathrm{K}}}$ .

\textit{Correct answer: } 699.6142

\item \hfill \textit{/ 10 pts}

A train traveling at 20 m/s is approaching a person on the platform, the engineer blowing its horn at 1,056 Hz. What frequency in Hz does the person on the platform hear? The speed of sound is 347 m/s.

\textit{Correct answer: } 1,120.5872

\item \hfill \textit{/ 10 pts}

A motorcyclist is traveling at 36 m/s away from a police car at rest. If the policeman blows his horn at a frequency of 685 Hz, what frequency in Hz does the motorcyclist hear? The speed of sound is 344 m/s.

\textit{Correct answer: } 613.314

\item \hfill \textit{/ 10 pts}

A car is traveling 10 m/s due east relative the ground. An ambulance is traveling on the same road, also going east at 43 m/s and away from the car, having already passed it. If the ambulance siren sounds at 1,130 Hz, what frequency does the driver of the car hear? The speed of sound in air: 343 m/s

\textit{Correct answer: } 1,033.3938

\item \hfill \textit{/ 10 pts}

A car is traveling due north at 25.8 m/s. A bicyclist behind the car is traveling at 12.8 m/s in the same direction. The driver of the car blows her horn at 916 Hz. What frequency does the bicyclist hear? Sound speed in air: 344 m/s.

\textit{Correct answer: } 883.7988

\item \hfill \textit{/ 10 pts}

A sound source creates a sound intensity at a certain distance. If the sound intensity is increased by a factor of 19.3, by how many decibels is the sound intensity level increased?

\textit{Correct answer: } 12.8556

\item \hfill \textit{/ 10 pts}

As the temperature rises, the fundamental frequency of wind instruments increases
, whereas as the temperature decreases, the frequency of wind instruments decreases
.

\textit{Answer 1:}
\begin{itemize}
  \item remains the same
  \item decreases
  \item \textbf{[correct]} increases
\end{itemize}
\textit{Answer 2:}
\begin{itemize}
  \item remains unchanged
  \item increases
  \item \textbf{[correct]} decreases
\end{itemize}

\end{enumerate}

\end{document}
